
\section{Generalisations}
There are two prominent ways to formalise and generalise the concept of derivations and their algebras.
\subsection{Lie Algebras}
The first way of formalisations are Lie algebras. 
\begin{defi}

An $R$ algebra $(\lieg,\mu)$ is called a Lie algebra if $\mu : \lieg \otimes \lieg \longrightarrow \lieg$ fulfills:
\bn
\item Jacobian identity:
$$\sum_{i=0}^2 \mu \circ (\mu \otimes id_{\lieg}) \circ \sigma_3^i(g_1 \otimes g_2 \otimes g_3) = 0,\ \mathrm{for~all}\ g_1 \otimes g_2 \otimes g_3 \in \lieg^{\otimes 3}.$$
Here, $\sigma_3 := (1,2,3) \in \mathrm{Gl}_R(\lieg^{\otimes3})$ is a $3$-cycle with $\sigma_3^3 = id$.
\item Antissymmetry:
$$\mu \circ \tau = - \mu,\ \mathrm{for}\ \tau = [g \otimes h \longmapsto h \otimes g].$$
\en
\end{defi}
\begin{koro}
Given any associative algebra $(A,\mu)$:
\bn
\item The commutator
$$[,] = [a \otimes b \longmapsto \mu(a \otimes b) - \mu(b \otimes a)] = \mu - \mu \tau,$$
defines a Lie algebra $(\lieg(A),[,])$ on $A$. As $R$ modules, they are identical.
\item The module of $R$ derivations on $A$ forms a Lie algebra $\mathrm{Der}_R(A)$ with commutator 
$$[,] : \lieg(\mathrm{End}_R(A)) \otimes \lieg(\mathrm{End}_R(A)) \longrightarrow \lieg(\mathrm{End}_R(A)),$$
restricted to the sub Lie algebra.
\en
\end{koro}
\bmk The proof is left to the reader. We shall extend our
\begin{defi}
Let $(\lieg,\mu)$ be an $R$-Lie algebra. We call an $R$-submodule $\lieh \subset \lieg$ an $R$-Lie subalgebra if
$$(\lieh,\mu\mid_{\frk{h}}) \subset (\lieg,\mu)$$
is a Lie algebra in its own right. Given two Lie algebras $(\lieg,\mu_{\lieg})$ and $(\lieh,\mu_{\lieh})$ and a homomorphism $f \in \mathrm{Hom}_R(\lieg,\lieh)$. We call $f$ a homomorphism of Lie algebras if
$$\xymatrix{
\lieg^{\otimes 2}\ar[r]^{f\otimes f}\ar[d]_{\mu_{\lieg}}&\lieh^{\otimes 2}\ar[d]^{\mu_{\lieh}}\\
\lieg \ar[r]_f& \lieh\\
}$$
commutes. We call a Lie subalgebra $\liea \subset \lieg$ a Lie ideal if
$$\mu(\liea \otimes_R \lieg) \subset \liea.$$
We call a Lie algebra $\lieg$ nilpotent if each sequence $(g_n) \in \lieg^{\nz}$ with
$$\mu^{n-1}(g_1 \otimes g_2 \otimes \ldots \otimes g_n) = 0$$
has finite length. Here, $\mu^n := \mu^{n-1}(\mu \otimes id_{\lieg^{\otimes (n - 1)}})$, $\mu^0 = id_\lieg$ and $\mu^1 = \mu$. For each Lie algebra $\lieg$ we get a chain of descending ideals: the so called $n$-th derived Lie algebra:
$$\mathcal{D}^n(\lieg) := \mu(\mathcal{D}^{n-1}(\lieg) \otimes \lieg),\ \mathrm{and}\ \mathcal{D}^0(\lieg) = \lieg.$$
We call $\lieg$ solvable if $\mathcal{D}(\lieg) = \mathcal{D}^1(\lieg)$ is nilpotent. 
\end{defi}

\bsp Let $R$ be a ring and $A = R[X]$ - its ring of polynomials. The endomorphism $\partial_X = [X \longmapsto 1] \in \mathrm{End}_R(A)$ is the classical example of a non-trivial $R$-derivation. Moreover, $\mathrm{Der}_R(R[X])$ contains $R[X]$-left modules:\\
\indent Claim: for any derivation $\partial : R[X] \longrightarrow R[X]$, the $R$-submodule
$$R[X].\partial \subset \der{R[X]}$$
is an $R[X]$-left module in $\der{R[X]}$.\\
\paragraph{Proof} Given a derivation $\partial \in \der{R[X]}$ and two polynomials $p, q \in R[X]$:
$$\bao{rcl}
[p \partial,q \partial] &=& p \partial(q \partial) - q \partial(p \partial)\\&&\\ &=& p \partial(q) \partial + p q \partial^2 - q \partial(p) \partial - q p \partial^2\\
&&\\
&=& (p \partial(q) - q \partial(p)) \partial + (p q - p q) \partial^2\\
&&\\
&=& (p \partial(q) - q \partial(p)) \partial \in R[X].\partial\\
\ea$$

Before we continue we will set for any $R$ module $M$ its Lie algebra to be $$(\mathfrak{gl}(M),\mu) = (\rendo{M}, [,])$$
the Lie algebra of R linear maps on $M$. In particular, if $M$ is free of rank $n$ than we use
$$\mathfrak{gl}_n(R)$$
synonymously.
\begin{defi}
Let $A$ be an $R$ algebra and $M$ an $R$ module. An $R$-linear map
$$\rho : A \longrightarrow \mathrm{End}_R(M)$$
is called a representation of $A$ if $\rho \in \mathrm{Alg}_R(A,\mathrm{End}_R(M))$ (i.e. a morphism of $R$ algebras). It is denoted by $(A,\rho)$.\\
Similarly, given a Lie algebra $\mathfrak{g}$ and an $R$ module $M$
we call
$$\rho : \mathfrak{g} \longrightarrow \mathfrak{gl}(M)$$
a representation of $\mathfrak{g}$ if $\rho \in \mathrm{Lie}(\mathfrak{g},\mathfrak{gl}(M))$ (i.e. a morphism of Lie algebras).
\end{defi}
For any Lie algebra $\mathfrak{g}$ the adjunct map
$$\bao{rrcl}
\mathrm{ad} : &\mathfrak{g}& \longrightarrow& \mathrm{End}_R(\mathfrak{g})
\\
\\
& x &\longmapsto& [x,] := [y \longmapsto [x,y]]
\ea$$
is an example of a non-trivial representation. 
\bmk We can generalize even further:\\
if $A$ is an associative (not necessarily commutative) $R$ algebra and $\mathcal{Z}(A)$ its center then 
$\dera$ is a left $\mathcal{Z}(A)$ module. If $A$ is also commutative then $\mathcal{Z}(A) = A$.
\paragraph{Proof} Given $a, b \in A$, $z \in \mathcal{Z}(A)$ and $\partial \ \dera$:
$$\bao{rcl}
z \partial (a b) &=& z (\partial(a) b + a \partial(b))\\
&&\\
&=& z \partial(a) b + z a \partial(b)\\
&&\\
&=& z \partial(a) b + a z \partial(b)\\
&\Leftrightarrow&\\
& & z \in \ker ad (a)\;\forall a \in A\\
\ea$$
~\\
\paragraph{Example} Fix $R$ as before and set $A := M_n(R)$, the $n \times n$ matrices over $R$ and $\mathfrak{gl}_n(R)$
\begin{description}
\item[Nilpotent Lie Algebra] is the subset of all strictly upper matrices:
$$\mathfrak{n}_n(R) := \left\{\left(a_{ij}\right) : a_{ij} = 0 \forall i \geq j\right\}.$$
This is an example of a Lie algebra of nilpotency index $\leq n$. The same applies to the strictly lower matrices 
\item[Solvable Lie Algebra] is the subset of all upper (lower) matrices:
$$\mathfrak{b}_n(R) := \left\{\left(a_{ij}\right) : a_{ij} = 0 \forall i > j\right\}.$$
\end{description}

\subsection{The Special Lie Algebra} 

Fix $R$ as before and set $A = R[X,Y]$. There are the two canonical derivations $\partial_X, \partial_Y \in \dera$. Now, set
$$\bao{rclcl}
D_X &=& \left[\bao{rcl}
X^i &\longmapsto& i X^{i-1}\\
Y^j &\longmapsto& -j Y^{j+1}\\
\ea\right] &=& \partial_X - Y^2 \partial_Y\\
&&&&\\
D_Y &=& \left[\bao{rcl}
X^i &\longmapsto& -i X^{i+1}\\
Y^j &\longmapsto& j Y^{j-1}\\
\ea\right] &=& \partial_Y - X^2 \partial_X\\
\ea$$
Computing the commutator of both derivations:
$$\bao{rcl}
[D_X,D_Y] &=& D_X D_Y - D_Y D_X\\
&&\\
 &=& (\partial_X - Y^2 \partial_Y)(\partial_Y - X^2 \partial_X)\\
 &&\\
 && - (\partial_Y - X^2 \partial_X) (\partial_X - Y^2 \partial_Y)\\
 &&\\
 &=& \partial_X \partial_Y - 2 X \partial_X - X^2 \partial_X^2\\
 &&\\
 && - Y^2 \partial_Y^2 + Y^2 X^2 \partial_Y \partial_X\\
 &&\\
 && - \partial_Y \partial_X + 2 Y \partial_Y + Y^2 \partial_Y^2\\
 &&\\ 
 && + X^2 \partial_X^2 - X^2 Y^2 \partial_X \partial_Y\\
 &&\\
 &=& \partial_X \partial_Y - 2 X \partial_X + X^2 Y^2 \partial_Y \partial_X \\
 &&\\
 && - \partial_Y \partial_X + 2 Y \partial_Y - X^2 Y^2 \partial_X \partial_Y\\
 &&\\
 &=& (1 - X^2 Y^2) [\partial_X,\partial_Y] + 2 Y \partial_Y - 2 X \partial_X\\
\ea$$
As the commutator $[\partial_X,\partial_Y]$ is trivial we arrive at
$$[D_X,D_Y] = 2 Y \partial_Y - 2 X \partial_X =: D_H.$$
Completing our computation for $[D_H,D_X]$ and $[D_H,D_Y]$, respectively:
$$\bao{rcl}
[D_H,D_X] &=& D_H D_X - D_X D_H\\
&&\\
&=& (2 Y \partial_Y - 2 X \partial_X)(\partial_X - Y^2 \partial_Y)\\
&&\\
&& - (\partial_X - Y^2 \partial_Y) (2 Y \partial_Y - 2 X \partial_X)\\
&&\\
&=& 2 Y \partial_Y \partial_X - 4 Y^2 \partial_Y - 2 Y^3 \partial_Y^2\\
&&\\
&& -2 X \partial_X^2 + 2 X Y^2 \partial_X \partial_Y\\
&&\\
&& - 2 Y \partial_X \partial_Y + 2 \partial_X + 2 X \partial_X^2\\
&&\\
&& + 2 Y^2 \partial_Y + 2 Y^3 \partial_Y^2 - 2 X Y^2 \partial_Y \partial_X\\ 
\ea$$
$$\bao{rcl}
[D_H,D_X] &=& 2 X Y^2 [\partial_X,\partial_Y] - 2 Y [\partial_X,\partial_Y]\\
&&\\
&& + 2 \partial_X - 2 Y^2 \partial_Y\\
&&\\
&=& 2 Y (X Y - 1) \underbrace{[\partial_X,\partial_Y]}_{0} + 2 \partial_X - 2 Y^2 \partial_Y\\
&&\\
&=& 2 (\partial_X - Y^2 \partial_Y)\\
&&\\
&=& 2 D_X\\
\ea$$
Considering the following map:
$$\tau : A \longrightarrow A,\;\left[\bao{rcl}
X^i &\longmapsto& Y^i\\
Y^j &\longmapsto& X^j\\
\ea\right]$$
which is an algebra automorphism of order two (i.e. $\tau^2 = id_A$). Then we get:
$$\bao{rclcl}
\tau \partial_X \tau &=& \left[\bao{rcl}
X^i &\longmapsto& 0\\
Y^j &\longmapsto& j Y^{j-1}\\
\ea\right] &=& \partial_Y\\
&&&&\\
\tau \partial_Y \tau &=& \tau^2 \partial_X \tau^2 &=& \partial_X\\
\ea$$
showing that both derivations are conjugated via $\tau$. Furthermore,
$$\bao{rcl}
\tau D_X \tau &=& \tau (\partial_X - Y^2 \partial_Y) \tau\\
&&\\
&=& \tau \partial_X \tau - \tau Y^2 \partial_Y \tau\\
&&\\
&=& \partial_Y - X^2 \tau \partial_Y \tau\\
&&\\
&=& \partial_Y - X^2 \partial_X\\
&&\\
&=& D_Y\\
&&\\
\tau D_H \tau &=& \tau (D_X D_Y - D_Y D_X) \tau \\
&&\\
&=& \tau D_X \tau^2 D_Y \tau - \tau D_Y \tau^2 D_X \tau\\
&&\\
&=& D_Y D_X - D_X D_Y\\
&&\\
&=& - D_H
\ea$$
giving us the additional identity $\tau D_Y \tau = D_X$. Equipped with these identities we get:
$$\bao{rclcl}
[D_H,D_Y] &=& D_H D_Y - D_Y D_H\\
&&\\
&=& D_H \tau D_X \tau - \tau D_X \tau D_H
&=& - \tau D_H D_X \tau + \tau D_X D_H \tau\\
&&\\
&=& - \tau [D_H,D_X] \tau
&=& -2 \tau D_X \tau\\
&&\\
&=& -2 D_Y\\
\ea$$
and the finding that our example algebra is an $\mathfrak{sl}_2(R)$ module (algebra). Lastly, consider the ideal
$$ I = \left<X Y - 1\right> \subset A.$$
Our generators $D_X$ and $D_Y$ give us the identities:
$$\bao{rcl}
D_Z\left( p (X Y - 1)\right) &=&  D_Z\left(p (X Y - 1)\right)\\
&&\\
&=& \underbrace{D_Z(p) (X Y - 1)}_{\in I} + p D_Z(X Y - 1)\\
&&\\
D_X(X Y - 1) &=& D_X(X) Y + X D_X(Y)\\
&&\\
&=& Y - X Y^2\\
&&\\
&=& -Y (X Y - 1) \in I\\
\\
D_Y(X Y - 1) &=& \tau D_X \tau(X Y - 1)\\
&&\\
&=& \tau D_X(X Y - 1)\\
&&\\
&=& -2 \tau (Y(X Y - 1)) = -2 X (X Y - 1) \in I\\
\ea$$
proving that $D_Z(I) \subset I$ for $Z = X, Y$. In addition, since both images are in $I$ so is $D_H(I)$ showing $I$ to be a differential submodule wrt. our new three derivations. Therefore, the projection
$$\pi : A \longrightarrow A/I = R[X,Y]/\left<X Y - 1\right> \simeq R[x,x^{-1}]$$
is a homomorphism of differential algebras and $A/I$ is an $\mathfrak{sl}_2(A/I)$ module algebra.
\begin{defi}

\end{defi}
\subsection{Coalgbras}
Speaking in a categorical manner, for every ring $R$ algebras are the dual category of the category of $R$ coalgebras. We compare the diagrams defining algebras and coalgebras. Again, $(A, \mu, \eta)$ is an unital associative $R$ algebra - we call an $R$ module $C$ a coalgebra if there is an $R$ linear map $\Delta : C \longrightarrow C^{\otimes2}$. In addition, we call a coalgebra $(C, \Delta)$ coassociative or counital for a given $\eps \in \mathrm{Hom}(C,R) = C^\ast$ if 
$$(\eps \otimes id_C) \Delta = (id_C \otimes \eps) \Delta = id_C$$
and the following diagrams commute:
\begin{longtable}{|c|cc|}
\hline
Coalgebras: &$$
\xymatrix{
C^{\otimes 3}  & C^{\otimes 2}\ar[l]_{id_C \otimes \Delta}\\
C^{\otimes 2} \ar[u]^{\Delta \otimes id_C}&C\ar[l]^{\Delta} \ar[u]_{\Delta}\\
}$$
&$$
\xymatrix{
&C^{\otimes 2}\ar[ld]_{\eps \otimes id_C}\ar[rd]^{id_C \otimes \eps}&\\
R \otimes C  &C\ar[u]_\Delta\ar[r]_{\simeq}\ar[l]^\simeq&C \otimes R \\
}$$\\
&Coassociativity & Counitality\\
\hline
&&\\
Algebras: &$$
\xymatrix{
A^{\otimes 3} \ar[r]^{\mu \otimes id_A} \ar[d]_{id_A \otimes \mu} & A^{\otimes 2}\ar[d]^{\mu}\\
A^{\otimes 2} \ar[r]_{\mu} &A\\
}$$
&$$
\xymatrix{
&A^{\otimes 2}\ar[d]^\mu&\\
R \otimes A \ar[r]_{\simeq}\ar[ru]^{\eta \otimes id_A} &A&A \otimes R \ar[lu]_{id_A \otimes \eta}\ar[l]^\simeq\\
}$$\\
&Associativity & Unitality\\
\hline
\end{longtable}
Clearly, each column is simply the inversion of arrows within the respective diagram. More interesting is the fact that the dual module $C^\ast$ for each coalgebra $(C, \Delta, \eps)$ is an algebra with multiplication
$$\mu_{C^\ast} := \Delta^\ast = \left[\alpha \otimes \beta \longmapsto (\alpha \otimes \beta) \Delta = \left[c \longmapsto \mu_R(\alpha \otimes \beta)\Delta(c) = \sum_{(c)} \alpha(c_{(1)}) \beta(c_{(2)})\right]\right],$$
with Sweedler notation $\Delta(c) = \sum_{(c)} c_{(1)} \otimes c_{(2)}$ and $\mu_R$ denoting the multiplication on $R$. With the dual of the counit we get a unit: $\eta_{C^\ast} = \eps^\ast = [1_R \longmapsto \eps]$.\\
We call a coalgebra $C$ cocommutative if
$$\xymatrix{
&C \ar[rd]^\Delta \ar[dl]_\Delta&\\
C^{\otimes2} \ar[rr]_\tau &&C^{\otimes2}\\
}$$
commutes for $\tau : C^{\otimes2} \longrightarrow C^{\otimes2}, c \otimes c' \longmapsto c' \otimes c$ the flip isomorphism.
\subsubsection{Group-likes and skew primitives}
Unless stated otherwise, we will always assume a coalgebra $C$ to be coassociative and counital.
\begin{defi}
We call an element $c \in C$ group-like if
$$\Delta(c) = c \otimes c.$$
We call an element $c$ $(g,h)$ skew primitive, for two group-like elements $g, h \in C$, if
$$\Delta(c) = g \otimes c + c \otimes h.$$
\end{defi}
Firstly, a group-like element $c \in C$ has $\eps(c) = 1$ as $(id_C \otimes \eps)\Delta(c) = id(c) \otimes \eps(c) = \eps(c) c \stackrel{!}{=} 1 \Leftrightarrow \eps(c) = 1$. Secondly, for two group-like $g, h \in C$ we have that a $(g,h)$ skew primitive element $c \in C$:
$$\bao{rcl}
c &\stackrel{!}{=}& (id_C \otimes \eps)\Delta(c) = id_C(g) \otimes \eps(c) + id_C(c) \otimes \eps(h)\\
&&\\
&\stackrel{!}{=}& (\eps \otimes id_C)\Delta(c) = \eps(g) \otimes id_C(c) + \eps(c) \otimes id_C(h)\\
\Leftrightarrow\\
\eps(c) &=& 0\\
\ea$$
\bmk Skew primitive elements are a generalisation of derivations - each derivation is simply an $(id_A,id_A)$ skew-primitive element. Another example are so called Ore extensions which are algebra extensions with respect to an $(\alpha,id_A)$ derivation $\partial_\alpha$ for some algebra (mono/auto) morphism $\alpha$:
$$B := A[X,\alpha,\partial_{\alpha}] = T(A \otimes X \otimes A)/\left<1_A \otimes X \otimes a - \alpha(a) \otimes X \otimes 1_A - \partial_{\alpha}(a) : a \in A\right>.$$
These types of algebra extensions are discussed studying quantum groups and quantum algebras and are beyond the scope of this paper.
\subsection{Bialgebras}
Given an $R$ algebra $D$ which is both and $R$ algebra and coalgebra:
\begin{defi}
then we call $D$ a bialgebra if $\mu, \eta$ are coalgebra homomorphisms and $\Delta, \eps$ are algebra homomorphisms:
$$\xymatrix{
D \ar[d] \ar[r] & D^{\otimes 2}\\
D & 
}$$
\end{defi}